\documentclass{article}
\usepackage[utf8]{inputenc}
\usepackage{amsmath,amssymb,booktabs,url,listings,xcolor}
\lstset{basicstyle=\ttfamily\small,breaklines=true}
\title{A Machine-Checked Soundness Proof for a DQBF Proof Checker\\
{\large Formally Verified Checking of DQRAT$+D^{\forall}$pure Certificates in Lean~4}}
% TODO(authors): co-authorship with the format/tool authors to be discussed.
\author{Zar Goertzel}
\date{June 2026 — draft}
\begin{document}
\maketitle

\begin{abstract}
Proof checkers are the trust anchors of automated reasoning: solvers may be
arbitrarily complex provided their certificates are validated independently.
That argument is only as strong as the checker itself. We present the first
machine-checked soundness proof of a proof checker for \emph{dependency
quantified Boolean formulas} (DQBF), covering the recent
DQRAT + $D^{\forall}$pure certificate format \cite{dqratcheck}. The proof, in
Lean~4~\cite{MouraUllrich2021}, is stated about the literal executable: if the
checker prints \texttt{s VERIFIED} --- on either of its two code paths --- the
input formula is semantically false. The development is \texttt{sorry}-free
and the top-level theorems depend only on the three standard classical axioms.
The mathematically hardest component is, to our knowledge, the first
formalization of \emph{full exhibition} of the reflexive resolution-path
dependency scheme~\cite{BeyersdorffBlinkhornChewSchmidtSuda2019,
WimmerWimmerSchollBecker2016}, which justifies the format's dependency-deletion
rule. Differential testing against the reference C++ implementation uncovered
two previously unknown defects in the latter (a garbage-collection bug that
both rejects valid proofs and crashes, and a parser non-termination), while one
semantic correction flowed in the opposite direction --- evidence that
verification pays as a development instrument, not only as a stamp. The
development was carried out largely by LLM agents under human direction; we
argue the resulting trust model is exactly the classical one: the Lean kernel,
not the proof's author, is the arbiter.
\end{abstract}

\section{Introduction}
The certificate paradigm dominates trustworthy automated reasoning: a solver
emits a proof which an independent, simpler checker validates
\cite{WetzlerHeuleHunt2014}. For propositional SAT the community went further
and verified the checkers themselves
\cite{CruzFilipeEtAl2017,Lammich2017}, so that trust reduces to a proof
assistant kernel. For \emph{quantified} and especially \emph{dependency
quantified} Boolean formulas (DQBF) --- where existential variables depend on
explicitly specified subsets of the universal variables, making satisfiability
\textsc{NEXP}-complete \cite{PetersonReif1979} --- proof formats are younger
\cite{HeuleSeidlBiere2014,dqratcheck} and no verified checker existed.

We verify the checker for DQRAT$+D^{\forall}$pure \cite{dqratcheck}, a format
whose proof rules include clause deletion, universal reduction, RUP/RAT-style
clause addition with dependency awareness (DQRATE/DQRATU), and --- its most
distinctive rule --- \emph{dependency deletion}: removing a universal $u$ from
the dependency set of an existential $x$ when a reachability analysis certifies
that no ``$u$-pure resolution paths'' connect the two polarities of $u$ to $x$.
Soundness of that rule is a nontrivial theorem of the dependency-scheme
literature \cite{WimmerWimmerSchollBecker2016,
BeyersdorffBlinkhornChewSchmidtSuda2019,VanGelder2011,SlivovskySzeider2012}.

\paragraph{Contributions.}
\begin{enumerate}
\item An executable DQRAT$+D^{\forall}$pure checker written in Lean~4 together
  with a machine-checked, end-to-end soundness proof: parser, every proof
  rule, and the command-line verdict are covered
  (\S\ref{sec:checker}). The claim-bearing theorems are
  \texttt{processProof\_sound'} and \texttt{parseDQDIMACS\_none\_sound}; the
  axiom audit reports only \texttt{propext}, \texttt{Classical.choice},
  \texttt{Quot.sound}.
\item The first formalization of full exhibition of the reflexive
  resolution-path dependency scheme, transcribed from
  \cite{BeyersdorffBlinkhornChewSchmidtSuda2019} with a
  definition-by-definition alignment table in the source
  (\S\ref{sec:deletion}); and a negative lesson from a failed first proof
  architecture (\S\ref{sec:descent}).
\item A differential-validation study against the reference C++ checker:
  verdict parity on the full upstream test suite, two new bugs found in the
  reference implementation, one conformance fix adopted on the verified side,
  and honest benchmarks (\S\ref{sec:differential}).
\item A report on carrying out such a development with LLM agents under
  kernel-checked ground truth, including where statement-level review by
  humans and adversarial models caught errors that type-checking cannot
  (\S\ref{sec:method}).
\end{enumerate}

\section{Background: DQBF and the certificate format}\label{sec:bg}
A DQBF prefix assigns each existential variable $x$ a dependency set
$D(x)$ of universal variables. Semantics is by Skolem functions: an
assignment $\mathit{sk}$ gives each $x$ a Boolean function of the values of
$D(x)$; the formula is \emph{true} iff some $\mathit{sk}$ satisfies the CNF
matrix under every universal assignment $\sigma$, and \emph{false} iff every
$\mathit{sk}$ is refuted by some $\sigma$. Our Lean development defines this
in ${\sim}176$ lines (\texttt{Semantics.lean}); this small file, together
with the statements of the two top theorems and the 38-line command-line
driver, is the entire semantic trust base.

A DQRAT refutation transforms the formula by truth-preserving steps ---
deletions (\texttt{d}), universal-variable additions (\texttt{a}),
dependency modifications (\texttt{e}), universal reduction lines
(\texttt{u}, with a DQRATU fallback gated by the $D^{\forall}$pure path
analysis), and RUP/DQRATE clause additions --- until an immediate
contradiction is derived; the original formula is then false.

\section{The verified checker}\label{sec:checker}
The artifact (${\sim}26$k lines of Lean across 23 modules) separates the
executable (parser, clause store with occurrence lists, unit propagation,
rule handlers) from its specification. The central invariant ladder
\texttt{Sound} $\subset$ \texttt{Correct} $\subset$ \texttt{FullCorrect}
captures, respectively, structural well-formedness; the \emph{action
boundary} (propagation quiescence, single trail level, and the property that
the current formula is a truth-preserving extension of the original); and
occurrence-list completeness. The last is not bureaucracy: the development
contains executable counterexamples (\texttt{Counterexamples.lean}) showing
that without it the DQRATE rule admits unsound acceptances --- the invariant
was \emph{forced} by a failed proof attempt, and \S\ref{sec:differential}
shows the reference implementation violating exactly this invariant after
garbage collection.

Every proof action is specified by a Hoare triple in the
\texttt{Std.Do} weakest-precondition framework, read as: \emph{if the
checker answers, the answer is right}; termination and completeness are
deliberately not claimed. The two top-level theorems cover the two
\texttt{s VERIFIED} print sites of the CLI:
\begin{quote}\small
\noindent\texttt{processProof\_sound'} says: if parsing the formula returns a
checker state \texttt{st} and the executable checker returns
\texttt{.Verified n} from that state, then the parsed formula is
semantically false.

\medskip
\noindent\texttt{parseDQDIMACS\_none\_sound} says: if
\texttt{parseDQDIMACS content = .ok none}, then there exist
\texttt{declaredMaxVar} and \texttt{st} such that the header's own
\texttt{maxVar} token yields \texttt{declaredMaxVar}, the inner parser run on
the tokenized comment-stripped input returns \texttt{.ok true st}, and the
resulting formula state is semantically false.
\end{quote}
The first assumes only the parse result and the checker verdict ---
\texttt{FullCorrect} is derived internally from the parser correctness
theorem, so there is no unverifiable side condition. The second covers
formulas refuted by unit propagation already during reading; its witness
state is pinned to the parser's own run on the input's token stream (a bare
existential would be vacuously true --- a statement-level subtlety caught in
review, \S\ref{sec:method}).

\section{Formalizing dependency deletion}\label{sec:deletion}
The \texttt{e}-line may delete $u$ from $D(x)$ when the checker's BFS
certifies the absence of cross-polarity $u$-pure resolution paths. Soundness
is \emph{full exhibition} of the reflexive resolution-path dependency
scheme: any model can be repaired into one whose Skolem functions ignore the
deleted dependency \cite{WimmerWimmerSchollBecker2016,
BeyersdorffBlinkhornChewSchmidtSuda2019,BeyersdorffBlinkhornPeitl2020}. We
transcribe the \emph{reformed models} proof of
\cite{BeyersdorffBlinkhornChewSchmidtSuda2019}: a two-stage transformation
copies each dependent existential's value from the $u$-flipped half-space
unless a resolution path \emph{refuses} the copy. The Lean module carries an
alignment table (paper Definitions 9--13, Lemmas 2 and 4, Theorem 8 $\to$
Lean identifiers); Lemma 3 is automatic in Skolem form and Lemma 5
unnecessary because the checker deletes one universal per step. A pleasant
proof-engineering fact: stating the refusal conditions as propositions
decided classically means only the \emph{completeness} direction of the BFS
specification is needed, never its soundness direction.

\subsection{A negative lesson: statement shape}\label{sec:descent}
Our first architecture attempted a witness-descent argument: a rank function
counting ``bad fibers'' and a restart obligation quantified over abstract
progress candidates. The obligation was eventually shown to admit
adversarial instances --- candidates scrambled off the deletion frontier
carrying no usable structure --- making the final lemma provably as hard as
the whole theorem. The descent was discarded (it is preserved in the
repository history) in favor of the literature's construction. The lesson we
draw: in formalization, \emph{finding the right statement} dominates proof
effort, and an abstraction that discards the provenance of its witnesses can
quietly convert a local lemma into the global theorem.

\section{Differential validation and benchmarks}\label{sec:differential}
We synchronized with the reference C++ checker at revision \texttt{d60d50e} and
compared verdicts on its full test suite, our regression corpus, and
generated instances (all scripts ship with the artifact). Results: exact
verdict parity on all nine upstream tests --- achieving which required
adopting one upstream semantic correction on our side (tautological input
clauses are kept, not dropped; both parser specification lemmas were
re-proved) --- and two previously unknown defects in the reference
implementation, found because the machine-checked \texttt{VERIFIED} of the
Lean checker is unimpeachable in a disagreement:
\begin{table}[t]
\centering
\small
\begin{tabular}{lcc}
\toprule
instance & upstream C++ \texttt{d60d50e} & Lean \\
\midrule
upstream tests (9 cases) & exact parity & exact parity \\
\texttt{proof\_missing\_zero} & HANG & FAILED (parse error) \\
\texttt{del5k} & FAILED (bug) & 0.00s VERIFIED \\
\texttt{del20k} & core dump & 3.0s VERIFIED \\
\bottomrule
\end{tabular}
\caption{Compact parity/benchmark slice from the artifact's differential study.}
\end{table}
(i) a garbage-collection bug in clause-deletion compaction that drops live
clauses from occurrence lists (valid proofs spuriously rejected from
${\sim}5000$ deletions; out-of-bounds crash at $20{,}000$); and
(ii) non-termination on a proof file whose final integer list lacks its
\texttt{0} terminator (a stream-extraction corner case: end-of-file reached
while skipping whitespace fails in the sentry, so the C++11
zero-on-failure rule never fires). Root-cause analyses and repro generators
are included; issue-ready upstream reports are included in the artifact. On performance: the verified
(non-watched) checker processes the benchmark family in $0.1$--$24$s
versus ${\sim}0.1$s for C++; an experimental watched-literals layer with
partial refinement proofs ships alongside, and completing its refinement to
the certified theorem is designated future work in the style of verified
watched-literal SAT solving \cite{BlanchetteFleuryWeidenbach2018}.

\section{Method: LLM agents under kernel ground truth}\label{sec:method}
The development was written predominantly by large-language-model agents
(several models, used adversarially against one another) under human
direction. We report this not as a novelty claim but because it sharpens the
trust question this paper answers: \emph{nothing} in the guarantee depends
on how the proofs were produced. The kernel checks every step; the axiom
audit bounds the ambient logic; the executable is the verified function.
What generation quality \emph{does} affect is the statements, and there
human and cross-model review earned its keep: an external model review
caught that the parse-time \texttt{VERIFIED} path was initially uncovered
(we proved it rather than scoping the claim down), and statement-level
review caught that the naive form of that theorem --- a bare existential
over states --- was vacuously true. Our working discipline: every claim in
this paper names a Lean identifier; sorry/axiom audits run in CI; diffs to
specifications get adversarial review; and counterexample modules document
why each load-bearing invariant exists.

\section{Related work}
Verified propositional certificate checking is mature
(LRAT \cite{CruzFilipeEtAl2017}, GRAT \cite{Lammich2017}); verified SAT
solving with watched literals \cite{BlanchetteFleuryWeidenbach2018} informs
our planned refinement. QBF preprocessing has QRAT
\cite{HeuleSeidlBiere2014}; dependency schemes originate with resolution
paths \cite{VanGelder2011,SlivovskySzeider2012}, were brought to DQBF by
\cite{WimmerWimmerSchollBecker2016}, reinterpreted and proved fully
exhibited in \cite{BeyersdorffBlinkhornChewSchmidtSuda2019}, and
strengthened in \cite{BeyersdorffBlinkhornPeitl2020}. To our knowledge no
prior proof checker for DQBF certificates — and no checker validating
dependency-scheme reasoning — has been formally verified.

\section{Conclusion}
A DQBF proof checker now exists whose \texttt{VERIFIED} verdict is a
theorem. The remaining gaps are stated plainly: completeness is not claimed;
performance parity awaits the watched-literals refinement. We hope the
artifact serves both as a trust anchor for DQBF solving (e.g., in
competition settings) and as a case study in LLM-assisted formalization
where the proof assistant, not the author, carries the trust.

\paragraph{Artifact.} Source, proofs, audits, benchmarks, and repro
generators: \url{https://github.com/godelclaw/dqrat-lean-checker}
(branch \texttt{codex/wip-parser-}%
\allowbreak\texttt{soundness-seam}).

\paragraph{Acknowledgments.} The development used Claude (Anthropic), Codex
(OpenAI) and GPT-5.x review models; and gratefully builds on the
DQRAT$+D^{\forall}$pure format and reference checker \cite{dqratcheck}.

\bibliographystyle{plain}
\bibliography{references}
\end{document}
